\section{Perspectiva de sistema: qué revisar primero cuando el rendimiento está limitado}
Para el entrenamiento de Agentic RL, el modo de inferencia y el uso de memoria de video no son dos temas independientes. El verdadero cuello de botella del sistema suele manifestarse así: por un lado se quiere aumentar el rendimiento del rollout, y por otro el context, el batch y la cache saturan la memoria de video. Entender la diferencia entre online y offline es, en esencia, entender cómo el sistema de entrenamiento negocia entre latencia y rendimiento.

\begin{knowledgebox}{Orden de revisión sugerido}
\begin{itemize}
  \item Primero revisa el número de tokens y el nivel de concurrencia de cada ronda de rollout
  \item Después revisa si la KV cache y el tamaño de lote saturan la memoria de video
  \item Por último decide si usar una orquestación de inferencia online u offline
\end{itemize}
\end{knowledgebox}

\subsection{Resumen de la sección}
El análisis de memoria de video no es un tema secundario, sino parte del diseño del sistema de RL. Solo cuando el rendimiento, la latencia y la memoria de video están equilibrados a la vez, la canalización de entrenamiento puede funcionar de manera estable.

\newpage
